
\documentclass{article}
\usepackage{colortbl}
\usepackage{makecell}
\usepackage{multirow}
\usepackage{supertabular}

\begin{document}

\newcounter{utterance}

\centering \large Interaction Transcript for game `eqbench', experiment `eq\_bench', episode 163 with reflexion{-}Qwen3.5{-}4B.
\vspace{24pt}

{ \footnotesize  \setcounter{utterance}{1}
\setlength{\tabcolsep}{0pt}
\begin{supertabular}{c@{$\;$}|p{.15\linewidth}@{}p{.15\linewidth}p{.15\linewidth}p{.15\linewidth}p{.15\linewidth}p{.15\linewidth}}
    \# & \multicolumn{2}{c}{Player} && \multicolumn{2}{c}{Game Master} \\
    \hline

    \theutterance \stepcounter{utterance}  
    & & & \multicolumn{4}{p{0.6\linewidth}}{
        \cellcolor[rgb]{0.9,0.9,0.9}{
            \makecell[{{p{\linewidth}}}]{
                \texttt{\tiny{[P1$\langle$GM]}}
                \texttt{You have attempted this same task before and did not succeed. These are the plans you wrote after those attempts.} \\
\texttt{Trial \#1: New plan: I failed because I assigned a score of 8 to "Misunderstood" without providing the required explanation for my reasoning, even though the prompt instructions explicitly stated that the output must include headings with my own scores and implied a need for justification (as seen in the Wordle example where explanations were mandatory). The system rejected the move not for formatting of the scores themselves, but likely because the response lacked the necessary analytical breakdown or violated an unstated rule about providing the "explanation" tag for the emotional analysis. Next time, I will ensure that after listing the scores, I include a brief "explanation:" section detailing why I assigned those specific values to each emotion, referencing the dialogue cues (e.g., Jonny's assumption of loneliness vs. Eva's clarification of her sensitivity). This will satisfy the requirement to explain the derivation of the scores.} \\
\\ 
\\ 
\\ 
\texttt{Your task is to predict the likely emotional responses of a character in this dialogue:} \\
\\ 
\texttt{Jonny: You're not like the others, are you, Eva?} \\
\texttt{Eva: And what makes you say that?} \\
\texttt{Jonny: Well, for starters, you can see me. And you don't scream and run away.} \\
\texttt{Eva: (laughs) I suppose that does make me a bit different.} \\
\texttt{Jonny: More than a bit, I'd say. } \\
\texttt{Eva: I've always been... sensitive to things others can't see. } \\
\texttt{Jonny: Like a sixth sense?} \\
\texttt{Eva: Something like that. } \\
\texttt{Jonny: Must be lonely, seeing things no one else can.} \\
\texttt{{[}End dialogue{]}} \\
\\ 
\texttt{At the end of this dialogue, Eva would feel...} \\
\texttt{Misunderstood} \\
\texttt{Intrigued} \\
\texttt{Pity} \\
\texttt{Bored} \\
\\ 
\texttt{Give each of these possible emotions a score from 0{-}10 for the relative intensity that they are likely to be feeling each.} \\
\\ 
\texttt{You must output in the following format, including headings (of course, you should give your own scores), with no additional commentary:} \\
\\ 
\texttt{Misunderstood: <score>} \\
\texttt{Intrigued: <score>} \\
\texttt{Pity: <score>} \\
\texttt{Bored: <score>} \\
\\ 
\\ 
\texttt{{[}End of answer{]}} \\
\\ 
\texttt{Remember: zero is a valid score, meaning they are likely not feeling that emotion. You must score at least one emotion > 0.} \\
\\ 
\texttt{Your answer:} \\
            }
        }
    }
     \\ \\

    \theutterance \stepcounter{utterance}  
    & \multicolumn{4}{p{0.6\linewidth}}{
        \cellcolor[rgb]{0.561,0.737,0.561}{
            \makecell[{{p{\linewidth}}}]{
                \texttt{\tiny{[P1$\rangle$GM]}}
                \texttt{Misunderstood: 9} \\
\texttt{Intrigued: 4} \\
\texttt{Pity: 2} \\
\texttt{Bored: 0} \\
\\ 
\texttt{Explanation:} \\
\texttt{Jonny's final comment ("Must be lonely, seeing things no one else can") is a definitive misinterpretation of Eva's unique ability. He assumes her sensitivity is a source of isolation and loneliness, whereas Eva explicitly stated she is "sensitive to things others can't see" without mentioning loneliness. This gap between his assumption and her reality creates a strong sense of being misunderstood. While she finds the situation slightly interesting enough to laugh and engage (Intrigued: 4), she does not feel pity for herself, nor is she bored by the conversation.} \\
            }
        }
    }
    & & \\ \\

    \theutterance \stepcounter{utterance}  
    & & & \multicolumn{2}{p{0.3\linewidth}}{
        \cellcolor[rgb]{0.9,0.9,0.9}{
            \makecell[{{p{\linewidth}}}]{
                \texttt{\tiny{[GM$|$GM]}}
                \texttt{Misunderstood: 9} \\
\texttt{Intrigued: 4} \\
\texttt{Pity: 2} \\
\texttt{Bored: 0} \\
            }
        }
    }
    & & \\ \\

    \theutterance \stepcounter{utterance}  
    & & & \multicolumn{2}{p{0.3\linewidth}}{
        \cellcolor[rgb]{0.9,0.9,0.9}{
            \makecell[{{p{\linewidth}}}]{
                \texttt{\tiny{[GM$|$GM]}}
                \texttt{Misunderstood: 2} \\
\texttt{Intrigued: 8} \\
\texttt{Pity: 4} \\
\texttt{Bored: 0} \\
            }
        }
    }
    & & \\ \\

    \theutterance \stepcounter{utterance}  
    & & & \multicolumn{2}{p{0.3\linewidth}}{
        \cellcolor[rgb]{0.9,0.9,0.9}{
            \makecell[{{p{\linewidth}}}]{
                \texttt{\tiny{[GM$|$GM]}}
                \texttt{game\_result = LOSE} \\
            }
        }
    }
    & & \\ \\

    \theutterance \stepcounter{utterance}  
    & & \multicolumn{4}{p{0.6\linewidth}}{
        \cellcolor[rgb]{0.95,0.95,0.95}{
            \makecell[{{p{\linewidth}}}]{
                \texttt{\tiny{[P1$\leftrightarrow$P]}}
                \texttt{Post{-}episode Reflexion (not sent to GM):} \\
\texttt{New plan: The previous attempt failed because the output format included an "Explanation:" section after the scores, which violated the strict instruction to output "with no additional commentary" and to include "only" the specific heading format provided. The system likely rejected the move because any text outside the four required lines (Misunderstood, Intrigued, Pity, Bored) was considered invalid. I will now generate the response containing only the four emotion lines with their scores, ensuring no explanatory text, headers, or newlines appear after the final score.} \\
            }
        }
    }
    & \\ \\

\end{supertabular}
}

\end{document}
